% Generated from ../chapters/25-a-broader-longevity-strategy.md
\chapter{A Broader Longevity Strategy}

The self-organizing-body framework should not replace conventional prevention or emerging biotechnology.

Cardiovascular risk, cancer screening, vaccination, dental care, infection treatment, hearing, vision, and metabolic disease require evidence-based attention. Catastrophic local failure can overwhelm global resilience.

A systems strategy has four layers.

\section{Layer 1: prevent avoidable damage}

Do not smoke. Maintain cardiorespiratory fitness, strength, sleep, nutrition, social connection, and appropriate medical care. Reduce exposure to known hazards.

\section{Layer 2: preserve adaptive reserve}

Train varied physical and cognitive capacities. Avoid narrowing life until small disturbances become crises.

\section{Layer 3: improve self-organization}

Use contemplative practice, slow movement, skilled training, and recovery to reduce maladaptive constraints and improve regulation.

\section{Layer 4: use targeted repair}

Drugs, surgery, devices, cell therapy, gene therapy, and future rejuvenation technologies may address barriers endogenous processes cannot cross.

The layers are complementary. A well-organized organism may respond better to therapy. A successful therapy may restore the capacity to practice and adapt.

\section{Beyond lifespan}

The objective is not merely additional years. It is preservation of agency: the ability to learn, contribute, move, relate, and reorganize.

Longevity without adaptive freedom is prolonged fragility. Healthspan is not only absence of disease; it is continued participation in life as a developing process.

\begin{center}\rule{0.5\linewidth}{0.5pt}\end{center}
